\begin{frame}{자주 하는 실수 2·4: GAN/Diffusion 판독}
  \begin{block}{실수 2 — ``생성자의 loss가 계속 줄어야 정상''이라고
  기대}
    GAN의 loss 곡선은 \textbf{등락} — $D$가 갱신될 때마다 $G$의
    손실 함수 \textbf{자체}가 바뀌므로, ``loss가 튀었다''가 반드시
    실패를 뜻하지 않음. (반대로 EM의 우도처럼 \textbf{단조증가}하는
    곡선을 기대하면 GAN의 ``정상''을 ``이상''으로 오독.)
    \; 실전에서는 \textbf{샘플 품질}을 직접 보고, $D$의 정확도가
    50\% 부근에서 \textbf{균형}을 유지하는지를 대시보드로.
  \end{block}
  \vskip0.4em
  \begin{block}{실수 4 — ``$D$의 정확도 50\%''를 ``잘 안 하고
  있다''로 읽기}
    \textbf{반대}다 — 50\%는 이 게임의 \textbf{목표} operating
    point(내쉬 균형). \textbf{정확도가 90\%+로 치솟는 쪽}이 진짜
    문제 신호 (위 ``실전 학습''의 실수 1 참고).
  \end{block}
\end{frame}
